% Generated draft: requires mathematical and source review.
\documentclass[11pt]{article}
\usepackage[margin=25mm]{geometry}
\usepackage{fontspec}
\setmainfont{Latin Modern Roman}
\usepackage{amsmath,amssymb,mathtools,unicode-math}
\setmathfont{Latin Modern Math}
\usepackage{enumitem,needspace,xurl,hyperref}
\setlength{\parindent}{0pt}
\setlength{\parskip}{5pt}
\setlength{\emergencystretch}{3em}
\newcommand{\N}{\mathbb{N}}
\newcommand{\Z}{\mathbb{Z}}
\newcommand{\Q}{\mathbb{Q}}
\newcommand{\R}{\mathbb{R}}
\newcommand{\C}{\mathbb{C}}
\newcommand{\F}{\mathbb{F}}
\newcommand{\A}{\mathbb{A}}
\newcommand{\PP}{\mathbb{P}}
\newcommand{\E}{\mathbb{E}}
\newcommand{\eps}{\varepsilon}
\newcommand{\dd}{\,\mathrm{d}}
\begin{document}
\section*{UnsolvedMath ID: 11000011}
\textbf{Source problem code:} catalogue entry
\subsection{Definitions and mathematical statement}
Let \(\operatorname{Mod}_g\) be the mapping class group of a closed oriented genus-\(g\) surface, with \(g\geq2\). A lattice in a Lie group \(G\) is a discrete subgroup with finite Haar-covolume; it is cocompact if \(\Gamma\backslash G\) is compact and noncocompact otherwise. No arithmeticity restriction is stated. The compound question asks:
\begin{enumerate}
\item Does there exist \(g\geq2\) for which \(\operatorname{Mod}_g\) contains a subgroup isomorphic to a cocompact lattice in \(SO(m,1)\) for some \(m\geq5\)? Likewise, does some \(\operatorname{Mod}_g\) contain a subgroup isomorphic to a noncocompact lattice in \(SO(m,1)\) for some \(m\geq4\)? The two dimension thresholds are different and retained.
\item Does some \(\operatorname{Mod}_g\), \(g\geq2\), contain a subgroup isomorphic to a cocompact lattice in \(SU(n,1)\), for some \(n\geq2\)?
\item For a fixed group \(\Gamma\) of any of those kinds and a fixed \(g\), must
\[\{H\leq\operatorname{Mod}_g:H\cong\Gamma\}/\text{conjugacy in }\operatorname{Mod}_g\]
be finite? This counts conjugacy classes of subgroups, not automatically all labelled embeddings modulo only target conjugacy.
\item Can a single fixed \(\operatorname{Mod}_g\) contain infinitely many pairwise nonisomorphic groups, each a cocompact lattice in \(SO(4,1)\)? The source singles out this lower real-hyperbolic dimension in light of an omitted example (3).
\end{enumerate}
The wording of the first two questions is existential in genus and dimension; no universal statement about every lattice or every genus is inserted.

The following discussion reports a finiteness assertion for representations \(\rho:\pi_1(V)\to\operatorname{Mod}_g\), with \(V\) a complex algebraic variety, admitting holomorphic representatives \(f:V\to\mathcal M_g\) in the indicated homotopy class. It attributes this to taking repeated hyperplane sections and then applying Arakelov--Parshin for curves, referring to Example 1.2 and Section 2.3. Its exact variety, nonconstant/isotrivial, orbifold and representation-equivalence conventions are absent; the coarse moduli space cannot silently be treated as a classifying space specifying a unique homotopy class from \(\rho\).

For representations without holomorphic representatives, the source proposes finding harmonic representatives and a Siu-type rigidity argument. It notes the difficulty that, in its \(\mathcal M_g\) setting, every loop can be freely homotoped outside every compact subset, and points to [DW] for then-recent progress. These are contextual reports and a proposed route, not additional hypotheses proving lattice embeddings or subgroup finiteness. The content of example (3), [DW] and their precise scope is missing.
\subsection{Short English statement}
Ask whether higher-dimensional real- or complex-hyperbolic lattices embed in mapping class groups, whether a fixed lattice has finitely many subgroup conjugacy classes, and whether one mapping class group contains infinitely many cocompact \(SO(4,1)\) lattice types. The holomorphic and harmonic-map discussion is qualified context.
\subsection{Sources}
\begin{itemize}
\item Local UnsolvedMath registry, permanent ID \(11000011\), exact source code \texttt{catalogue entry}.
\item Farb, \emph{Problems on Mapping Class Groups and Related Topics}, Question 2.5, PDF page 25: \url{https://www.math.uchicago.edu/~farb/papers/mcgbook.pdf}. Registry access 2026-07-29; pdf-font-structure extraction.
\item The source invokes example (3), Example 1.2, Section 2.3, Arakelov--Parshin, Siu-type rigidity and [DW], without reproducing the precise statements or the [DW] bibliography.
\item The research summary mentions Kahn--Markovic and Leininger--Reid in surface-subgroup context without precise citations; these claims are not used as evidence for the higher-dimensional lattice questions.
\end{itemize}
\end{document}
